\section{Metodología}
\label{sec:metodologia}

Esta sección explica, paso a paso, cómo obtiene el sistema sus resultados: qué hace con cada
fotograma desde que llega de una cámara hasta que se convierte en una alerta o en un reporte,
cómo se evalúa lo que produce y cómo se trabajó para llegar a él. Describe el comportamiento
del prototipo vigente; los requisitos que ese comportamiento satisface se especifican en la
Sección~\ref{sec:requerimientos} y su diseño en la Sección~\ref{sec:diseno-sistema}.

\subsection{Visión general}
\label{subsec:met-vision}

El sistema procesa el video por muestras, no cuadro a cuadro. En cada muestra lee un
fotograma de cada cámara, detecta a las personas, descarta lo que no son personas o queda
fuera del área de interés, sigue a cada persona dentro de su cámara, proyecta la posición de sus
pies sobre el plano del espacio, decide su identidad entre cámaras y, con las posiciones de las
personas confirmadas, calcula la ocupación de cada zona y detecta aglomeraciones. Las alertas,
la bitácora, el panel comercial y los reportes se alimentan de ese análisis
(Figura~\ref{fig:flujo-metodologia}). Antes de procesar video, el operador configura el espacio:
el plano, las cámaras, las zonas y la calibración de cada cámara.

\begin{figure}[H]
\centering
\resizebox{\anchofigura}{!}{%
\begin{tikzpicture}[
  etapa/.style={draw=black,thick,rounded corners=2pt,align=center,font=\small,
                minimum width=6.6cm,minimum height=1.15cm,inner sep=3pt},
  config/.style={draw=black,dashed,rounded corners=2pt,align=center,font=\small,
                 minimum width=4.4cm,minimum height=1.15cm,inner sep=3pt},
  flecha/.style={-{Stealth[length=2.2mm]},thick}]
  \node[etapa] (e1) at (0,0) {\textbf{Fuente de video}\\archivo, cámara USB o flujo de red};
  \node[etapa,below=0.5cm of e1] (e2) {\textbf{Detección}\\YOLO11, clase persona};
  \node[etapa,below=0.5cm of e2] (e3) {\textbf{Filtros}\\zona útil, tamaño, elementos estáticos};
  \node[etapa,below=0.5cm of e3] (e4) {\textbf{Seguimiento}\\punto de los pies, Kalman, IoU};
  \node[etapa,below=0.5cm of e4] (e5) {\textbf{Proyección al plano}\\homografía de cada cámara};
  \node[etapa,below=0.5cm of e5] (e6) {\textbf{Identidad}\\OSNet, entre oclusiones y cámaras};
  \node[etapa,below=0.5cm of e6] (e7) {\textbf{Análisis por zonas}\\ocupación, permanencia, aglomeración};
  \node[etapa,below=0.5cm of e7] (e8) {\textbf{Alertas y reportes}\\pantalla, correo, bitácora, panel, archivos};
  \foreach \a/\b in {e1/e2,e2/e3,e3/e4,e4/e5,e5/e6,e6/e7,e7/e8} { \draw[flecha] (\a) -- (\b); }
  \node[config,right=1.3cm of e3] (c1) {\textbf{Perfil de hardware}\\variante del detector,\\resolución y muestreo};
  \node[config,right=1.3cm of e5] (c2) {\textbf{Configuración del espacio}\\plano, cámaras, zonas\\y calibración};
  \draw[flecha,dashed] (c1.west) -- ++(-0.65,0) |- (e2.east);
  \draw[flecha,dashed] (c2.west) -- (e5.east);
  \draw[flecha,dashed] (c2.south) |- (e7.east);
\end{tikzpicture}}
\caption{Flujo de procesamiento del sistema, desde la fuente de video hasta las alertas y
los reportes. Los recuadros discontinuos son configuraciones previas que condicionan las
etapas.}
\label{fig:flujo-metodologia}
\end{figure}

\subsection{Etapas del procesamiento}
\label{subsec:met-etapas}

Cada etapa se describe por lo que recibe, lo que hace, lo que entrega y la razón por la que se
eligió esa solución.

\subsubsection{Adquisición}
\label{subsec:met-adquisicion}

\textbf{Entrada.} Un archivo de video grabado, una cámara USB o un flujo de red (RTSP o una
cámara IP), por cada cámara del proyecto. Una misma sesión no mezcla archivos y cámaras en
vivo, porque sus relojes no son equivalentes.

\textbf{Proceso.} Cada fuente se abre una sola vez y se decodifica en su propio hilo. Los
fotogramas de más de 1280 píxeles de ancho se reducen a ese ancho. Antes de empezar, el
sistema determina un perfil de hardware: si el equipo no tiene una GPU utilizable usa la
variante nano de YOLO11 con 960 píxeles de entrada; con GPU elige la variante nano, pequeña o
mediana según la memoria de video disponible (menos de 6 GB, de 6 a 16 GB, o 16 GB o más) y
resoluciones de 960 a 1920 píxeles. Cuando alguna cámara está a tres metros de altura o más,
sube la resolución, porque las personas ocupan pocos píxeles. El operador elige además un
modo de rendimiento: preciso (una muestra cada 0,2 segundos), equilibrado (cada 0,4 segundos y
800 píxeles) o rápido (cada 0,6 segundos y 640 píxeles); en modo automático y sin GPU, el
sistema usa el preciso hasta con dos cámaras, el equilibrado hasta con cuatro y el rápido con
más.

\textbf{Salida.} En cada paso, un fotograma por cámara activa, con su instante en la fuente.

\textbf{Por qué.} El prototipo debe funcionar en computadoras portátiles sin GPU dedicada.
Muestrear cada pocas décimas de segundo es suficiente para fenómenos que se forman en
minutos, como una aglomeración, y reduce a la vez el cómputo y el volumen de imágenes
tratadas.

\subsubsection{Configuración del espacio}
\label{subsec:met-configuracion}

\textbf{Entrada.} El plano del espacio (uno de los planos del aeropuerto, una imagen, un PDF o
un archivo DXF, o un dibujo hecho en la consola) con sus dimensiones en metros, y una imagen
de cada cámara.

\textbf{Proceso.} El operador ubica cada cámara sobre el plano, marca en su imagen la zona útil
que debe analizarse y, si la cámara ve una puerta, un acceso formado por un área exterior y otra
interior. Dibuja las zonas de interés y les asigna un tipo y, si corresponde, una regla de
aglomeración propia. Para calibrar, marca puntos del suelo visibles en el video (esquinas,
cruces de baldosas) y su posición correspondiente en el plano. Con cuatro puntos el sistema
calcula una homografía exacta; con cinco o más, la estima de forma robusta descartando los
puntos incoherentes y la reajusta con los restantes, y además comprueba cada punto dejándolo
fuera del ajuste. Rechaza las calibraciones con puntos repetidos o casi alineados. Por último, el
operador declara qué cámaras están relacionadas y si sus relojes están sincronizados.

\textbf{Salida.} La configuración del proyecto: plano, cámaras, zonas, reglas y una homografía
por cámara calibrada, junto con la región de la imagen en la que esa homografía es fiable.

\textbf{Por qué.} Las personas caminan sobre el suelo, y una homografía entre el suelo de la
imagen y el plano basta para pasar de píxeles a metros sin conocer los parámetros internos de
cada cámara. Pedir los puntos al operador, en lugar de estimarlos, hace la calibración
verificable y repetible.

\subsubsection{Detección}
\label{subsec:met-deteccion}

\textbf{Entrada.} Los fotogramas de todas las cámaras de la muestra.

\textbf{Proceso.} El detector YOLO11 (Jocher y Qiu, 2024), preentrenado sobre el conjunto
COCO, procesa los fotogramas de todas las cámaras en un solo lote y conserva únicamente la
clase persona. El umbral de confianza es de 0,25, y de 0,15 cuando hay cámaras elevadas. La
supresión de cajas duplicadas usa un umbral de superposición de 0,5, más estricto que el
habitual, para no contar dos veces a una misma persona dentro de un grupo.

\textbf{Salida.} Por cada persona detectada, una caja, su confianza y el punto de los pies, que
es el centro del borde inferior de la caja.

\textbf{Por qué.} YOLO entrega cajas completas en una sola pasada y con variantes de distinto
costo, y el borde inferior de la caja da el punto que la homografía puede proyectar al suelo
(Sección~\ref{subsec:yolo}).

\subsubsection{Filtros de falsos positivos}
\label{subsec:met-filtros}

\textbf{Entrada.} Las detecciones de cada cámara.

\textbf{Proceso.} Se aplican tres filtros. El primero descarta las detecciones cuyo punto de los
pies cae fuera de la zona útil de la imagen o fuera del área de trabajo del plano; la consola las
sigue dibujando en gris para que se vea que el detector sí las encontró. El segundo es un filtro
de tamaño relativo: cada cámara aprende la altura típica de las personas que detecta con
confianza alta (0,6 o más) y descarta las detecciones dudosas que miden menos del 62\,\% de esa
referencia, salvo que estén cortadas por el borde de la imagen. El tercero actúa después del
seguimiento y suprime los elementos estáticos: una trayectoria que en dos segundos no se ha
movido más de una cuarta parte del ancho de su caja y cuya confianza media es baja se trata
como mobiliario, un cartel o un reflejo. Una persona quieta pero bien detectada se conserva.
Los dos últimos filtros pueden desactivarse en la configuración.

\textbf{Salida.} Las detecciones que se consideran personas, y el número de descartes por
cámara.

\textbf{Por qué.} A los umbrales de confianza bajos que exigen las cámaras elevadas, el detector
confunde con personas objetos como gorros, carteles o maniquíes. Los dos filtros usan señales
de la propia escena, sin entrenar nada, y se adaptan a cámaras de distinta altura.

\subsubsection{Seguimiento}
\label{subsec:met-seguimiento}

\textbf{Entrada.} Las detecciones filtradas de cada cámara.

\textbf{Proceso.} Cada persona se sigue sobre el punto de sus pies con un filtro de Kalman de
velocidad constante que predice dónde estará en la muestra siguiente. La asociación entre
predicciones y detecciones se hace en dos etapas, como en ByteTrack (Zhang et al., 2022):
primero con las detecciones de confianza alta (0,4 o más) y después con las de confianza baja,
para no perder a una persona parcialmente oculta, con una puerta espacial de 60 píxeles. Dentro
de esa puerta, las coincidencias se ordenan además con dos señales al estilo de BoT-SORT
(Aharon et al., 2022): la superposición (IoU) entre la caja predicha y la detectada, y una firma
de color de la ropa del torso. Una persona que deja de detectarse se conserva durante 45
muestras antes de darse por perdida.

\textbf{Salida.} Por cámara, la lista de personas seguidas, cada una con un identificador local
estable, su caja, su punto de los pies y su velocidad.

\textbf{Por qué.} Seguir el punto de los pies, y no solo la caja, mantiene la geometría que la
homografía necesita. La segunda etapa de asociación y las señales de caja y color reducen los
cambios de identidad en grupos y oclusiones cortas.

\subsubsection{Proyección al plano}
\label{subsec:met-proyeccion}

\textbf{Entrada.} El punto de los pies de cada persona seguida y la homografía de su cámara.

\textbf{Proceso.} El punto se transforma con la homografía a coordenadas del plano, en metros
si el plano tiene escala. Solo se proyectan los puntos situados dentro de la región fiable de la
cámara, que es la zona de la imagen cubierta por los puntos de calibración, ampliada; fuera de
ella el sistema no extrapola: la persona sigue apareciendo en el video, pero no en el plano, y la
consola avisa al operador para que añada referencias en esa zona.

\textbf{Salida.} La posición de cada persona en el plano común a todas las cámaras.

\textbf{Por qué.} Un plano común permite sumar lo que ven varias cámaras, medir distancias en
metros para las reglas de aglomeración y usar la posición para decidir si dos observaciones
corresponden a la misma persona. No extrapolar evita los errores grandes de una homografía
lejos de sus puntos de referencia.

\subsubsection{Identidad}
\label{subsec:met-identidad}

\textbf{Entrada.} Las personas seguidas en todas las cámaras, con su recorte de imagen y, si la
cámara está calibrada, su posición en el plano.

\textbf{Proceso.} Para cada vista confiable de una persona (confianza suficiente, altura mínima,
lejos del borde y sin oclusión), el sistema calcula un vector de apariencia OSNet (Zhou et al.,
2019) sobre el recorte de cuerpo entero y lo acumula en el tramo de seguimiento al que
pertenece. Con esos vectores mantiene el identificador cuando una persona se pierde un
momento detrás de otra y vuelve a aparecer, y lo extiende a otra cámara cuando el parecido de
apariencia es suficiente y el tiempo y, con calibración, la posición son coherentes, siempre que
las cámaras estén relacionadas y sus relojes declarados sincronizados. Una persona recibe un
identificador provisional hasta reunir suficientes vistas y solo entonces uno confirmado; las
asociaciones ambiguas se marcan como inciertas. Con la memoria de apariencia activa, que es
la configuración por defecto, el sistema compara además cada identidad nueva con las
personas guardadas de sesiones anteriores dentro del plazo de retención. Esa memoria es la
causa del hallazgo no conforme de RNF-01 descrito en la Sección~\ref{subsec:privacidad} y en
el apartado B.9 del Anexo~\ref{anx:marco}. Al cerrar la sesión, los tramos se reagrupan y las
personas se renumeran de 1 a N.

\textbf{Salida.} Para cada observación, un identificador provisional o confirmado y su estado de
asociación: local, estimada, reidentificada o incierta.

\textbf{Por qué.} Sin identidad entre cámaras, una persona vista por dos cámaras se contaría dos
veces. La sola proximidad en el plano no basta cuando las cámaras no comparten campo de
visión o cuando la calibración es imperfecta, y por eso se añadió el descriptor de apariencia.
Contar solo los identificadores confirmados evita que una detección fugaz infle los conteos.

\subsubsection{Análisis}
\label{subsec:met-analisis}

\textbf{Entrada.} Las posiciones en el plano de las personas con identificador confirmado y la
configuración de zonas y reglas.

\textbf{Proceso.} En cada muestra el sistema calcula:
\begin{itemize}[leftmargin=2.2em,itemsep=0.1em]
  \item \textbf{Ocupación por zona:} cuántas personas hay dentro de cada polígono, cuántas
  personas-segundo acumula, su pico y cuántos identificadores distintos la visitaron, además de
  un mapa de calor del plano por celdas.
  \item \textbf{Aglomeraciones con la regla global:} alrededor de cada persona se traza un círculo
  del radio configurado; si dentro hay al menos la cantidad mínima de personas, es una
  concentración candidata, y se descartan los círculos que repiten a otro mayor. Cada
  concentración se empareja con la de la muestra anterior más cercana para medir cuánto lleva
  formada, y genera una alerta cuando esa duración alcanza el tiempo de permanencia
  configurado.
  \item \textbf{Aglomeraciones con reglas por zona:} una zona con regla propia alerta cuando
  contiene al menos su cantidad mínima de personas durante su tiempo de permanencia.
  \item \textbf{Flujo por los accesos:} el paso de personas entre el área exterior y la interior de
  cada acceso se cuenta como entrada o salida.
  \item \textbf{Permanencia por negocio:} a partir de la grabación de la sesión, la analítica
  comercial calcula visitas, permanencias, regresos y rutas frente a cada negocio del piso.
\end{itemize}

\textbf{Salida.} La ocupación y el historial de cada zona, las concentraciones con su duración y
su estado de alerta, los cruces por acceso y las métricas comerciales.

\textbf{Por qué.} Las tres condiciones de una aglomeración (cuántas personas, en cuánto espacio
y durante cuánto tiempo) son las variables que la Sección~\ref{sec:planteamiento} define como
objeto de medición. Exigir que la concentración se sostenga en el tiempo evita alertar por una
coincidencia de unos segundos, y las reglas por zona permiten que una cola tolere más
personas que una sala de espera.

\subsubsection{Salidas}
\label{subsec:met-salidas}

\textbf{Entrada.} Los resultados del análisis.

\textbf{Proceso.} Cada alerta se muestra en la consola y se registra como incidente en la
bitácora, donde el personal la marca como revisada, falsa alarma o resuelta. Si el aviso por
correo está activo, el sistema envía un correo por incidente nuevo, sin repetir una misma zona
antes de 120 segundos. El dashboard ejecutivo y el panel comercial presentan las zonas
ordenadas por ocupación y por oportunidad comercial. Los reportes se exportan en PDF, Excel,
CSV o JSON, y cada uno declara si sus datos proceden de una simulación, de videos de prueba o
de una fuente real sin validación de precisión. Cada sesión queda además grabada en el equipo,
con los identificadores y las posiciones de cada persona, para poder revisarla después.

\textbf{Salida.} Avisos para el personal, un registro auditable de incidentes y los archivos de
reporte.

\textbf{Por qué.} El sistema alerta y no actúa: la decisión queda siempre en manos de una
persona (medida M-10). La etiqueta de origen en cada reporte impide presentar como real un
resultado sintético. La grabación de la sesión, en cambio, conserva datos por persona sin un
plazo automático, y es la causa del hallazgo de la medida M-03 (apartado B.9 del
Anexo~\ref{anx:marco}).

\subsection{Evaluación de los resultados}
\label{subsec:met-evaluacion}

\textbf{Métricas.} La detección se evalúa con la exhaustividad (\emph{recall}), la fracción de
personas reales que el sistema encuentra, y la precisión, la fracción de detecciones que son
personas reales. Ambas se calculan emparejando las cajas del sistema con las de referencia
mediante asignación húngara y exigiendo una superposición (IoU) de al menos 0,5. El conteo se
evalúa con el error absoluto medio (MAE) entre las personas contadas y las reales por cuadro,
junto con su sesgo, que indica si el sistema cuenta de más o de menos. El rendimiento se mide
en cuadros procesados por segundo y en milisegundos por etapa: el propio sistema registra en
cada ciclo el tiempo del detector, del seguimiento y del descriptor de apariencia. La identidad
entre cámaras se evalúa con la métrica IDF1, la fracción de detecciones que conservan la
identidad correcta (Sección~\ref{subsec:variables-restricciones}).

\textbf{MOT20, solo para evaluación.} Para medir la detección y el conteo en escenas densas se
usa el conjunto público MOT20 (Dendorfer et al., 2020), cuya licencia no permite el uso
comercial.\verificar{licencia de MOT20: el equipo la registra como CC BY-NC-SA 3.0; contrastar
con los términos publicados en el sitio de MOTChallenge.} Por esa razón se emplea únicamente
para evaluar, nunca para entrenar ni para ajustar un modelo que se entregue a LAP. El protocolo
usa las secuencias MOT20-01 y MOT20-03, toma uno de cada cinco cuadros, que equivale al paso
de 0,2 segundos del modo preciso, considera como referencia a las personas visibles al menos en
un 25\,\% e ignora a las más ocultas, y mide el sistema en tres puntos: a la salida del detector,
tras los filtros y al final del seguimiento, que es lo que el sistema cuenta. Medir en tres puntos
permite atribuir cada pérdida a la etapa que la produce. La Tabla~\ref{tab:mot20} resume el
resultado con la configuración vigente.\verificar{experimento MOT20: realizado sobre una copia
de trabajo del prototipo que todavía no se ha integrado en su versión principal; confirmar las
cifras al integrarlo.}

\begin{table}[H]
\centering
\small
\caption{Evaluación con MOT20 de la configuración vigente, al final del seguimiento.}
\label{tab:mot20}
\begin{tabular}{@{}L{3.6cm}C{1.9cm}C{1.9cm}C{2.0cm}C{2.0cm}@{}}
\toprule
\textbf{Secuencia} & \textbf{Recall (\%)} & \textbf{Precisión (\%)} & \textbf{MAE (\%)} & \textbf{FPS} \\
\midrule
MOT20-01 (diurna) & 70,5 & 94,0 & 25,0 & 10,0 \\
MOT20-03 (nocturna, cámara alta) & 65,5 & 91,4 & 28,4 & 3,4 \\
\bottomrule
\end{tabular}
\end{table}

Las cifras se obtuvieron en un equipo con una GPU GTX 1650 de 4 GB, sin el descriptor de
apariencia, y el sistema cuenta de menos en ambas secuencias. En MOT20-01 el detector
encuentra al 82,8\,\% de las personas, pero el filtro de tamaño y el umbral de confianza del
seguimiento bajan esa cifra al 70,5\,\%. Ninguna de estas cifras predice el desempeño en el
Aeropuerto Jorge Chávez (RNF-05).

\textbf{Pruebas automatizadas.} El prototipo incluye 322 pruebas unitarias y de integración que
comprueban, entre otros puntos, la validación de la configuración, la calibración, el
seguimiento, el motor de identidad, la memoria de apariencia, el aislamiento entre proyectos y
la generación de reportes.\verificar{resultado de la ejecución completa de las 322 pruebas
sobre la versión entregada: no se ejecutaron al redactar esta sección.}

\textbf{Modo demostración.} Para comprobar la cadena de análisis y las salidas sin necesidad de
video, el sistema ofrece un modo de demostración que genera dieciocho personas sintéticas con
trayectorias periódicas sobre el plano y cruces simulados en los accesos. Cada persona
sintética lleva un identificador propio de la demostración y su asociación se marca como
sintética, y los paneles y reportes de esa sesión se rotulan como simulación (RNF-04). El modo
demostración verifica que el análisis, las alertas y los reportes funcionan; no mide la exactitud
del sistema.

\subsection{Organización del trabajo}
\label{subsec:met-trabajo}

El proyecto avanzó por iteraciones semanales registradas en actas de reunión, en las que se
asignaban las tareas de cada frente, y se presentó al curso en avances quincenales
(Sección~\ref{sec:planificacion}).\verificar{actas semanales: el repositorio del proyecto solo
conserva el acta de la semana 1, del 26 de agosto de 2026.} El detector evolucionó en tres
pasos. La primera versión usó un detector clásico basado en histogramas de gradientes
orientados, HOG (Dalal y Triggs, 2005), como referencia sin costo de GPU, que se abandonó por
su baja sensibilidad con personas parcialmente ocultas.\verificar{motivo del abandono de HOG:
no hay una medición documentada que lo respalde.} La segunda incorporó P2PNet (Song et al.,
2021), que localiza muy bien personas en multitudes, pero que se retiró por dos razones: localiza
cabezas y no pies, lo que impide proyectar la posición al plano sin un error no validado, y en
CPU tardaba del orden de 3,5 segundos por cada 0,2 segundos de video. La tercera, vigente,
adoptó YOLO11, que entrega el punto de los pies, funciona en CPU y escala con el hardware. La
identidad entre cámaras siguió un camino paralelo: empezó con la proximidad en el plano y un
histograma de color como desempate, y pasó a OSNet cuando la proximidad dejó de bastar.

La privacidad se trató desde el diseño y no al final: cada decisión técnica se contrastó con las
restricciones éticas y legales de la Sección~\ref{subsec:restricciones}, con el marco normativo
del Anexo~\ref{anx:marco} y con el análisis de impacto del Anexo~\ref{anx:impacto}. Ese
contraste es el que identificó los dos hallazgos que el equipo debe corregir: la memoria de
apariencia activa por defecto (RNF-01) y la persistencia de posiciones individuales en las
grabaciones (M-03).
